% ex: ts=2 sw=2 sts=2 et filetype=tex
% SPDX-License-Identifier: CC-BY-SA-4.0
\documentclass[aspectratio=169]{beamer}

\usetheme{Boadilla}

\setbeamertemplate{navigation symbols}{}
\graphicspath{{../img/}{../styles/iteso/}}
\usepackage{../styles/iteso/beamercolorthemeiteso}

\usepackage[utf8]{inputenc}
\usepackage{graphicx}
\usepackage{listings}

\title[Elementos básicos de un programa]{Estructuras secuenciales:\\
Elementos básicos de un programa}

% ex: ts=2 sw=2 sts=2 et filetype=tex
% SPDX-License-Identifier: CC-BY-SA-4.0

% Lenguaje PSeInt 
\lstdefinelanguage{PSeInt}{
  keywords={Algoritmo, FinAlgoritmo, Funcion, Escribir, Leer, Si, Entonces,
  SiNo, Fin, Segun, Hacer, De, Otro, Modo, Mientras, Repetir, Hasta, Que,
  Para, Hasta, Con, Paso, Sin, Saltar},
  sensitive=False,
  morecomment=[l]{//},
  morecomment=[s]{/*}{*/},
  morestring=[b]',
  morestring=[b]",
  alsoletter={<>=-:;},
  alsodigit={-}
}

\lstdefinestyle{pseudocodigo} {%
  % General design
  backgroundcolor=\color{light-gray},
  basicstyle=\small,
  numberstyle=\tiny\color{codeNumbers},
  xleftmargin={0.2cm},
  numbersep=5pt,
  numbers=left,
  stepnumber=1,
  firstnumber=1,
  numberfirstline=true,
  % Code design
  identifierstyle=\color{black},
  keywordstyle=\color{blue}\bfseries,
  keywordstyle={[2]\color{red}},
  stringstyle=\color{codeString}\ttfamily,
  commentstyle=\color{codeNumbers}\ttfamily,
  % Code
  language=PSeInt,
  alsodigit={.:;},
  tabsize=2,
  showtabs=false,
  showspaces=false,
  showstringspaces=false,
  extendedchars=true,
  breaklines=true,
  breakatwhitespace=false,
  captionpos=b,
  keepspaces=true,
  upquote=true,                      % requires textcomp
  % Español
  literate=%
  {á}{{\'a}}1 {é}{{\'e}}1 {í}{{\'i}}1 {ó}{{\'o}}1 {ú}{{\'u}}1
  {Á}{{\'A}}1 {É}{{\'E}}1 {Í}{{\'I}}1 {Ó}{{\'O}}1 {Ú}{{\'U}}1
  {¡}{{!`}}1  {¿}{{?`}}1
  {ñ}{{\~n}}1 {Ñ}{{\~N}}1
}


\author{Miguel Bernal Marin} % Your name
\institute[ITESO] % Your institution as it will appear on the bottom of every slide, may be shorthand to save space
{
 ITESO, Universidad \\
 Jesuita de Guadalajara \\% Your institution for the title page
\medskip
\textit{miguel.bernal@iteso.mx} % Your email address
}
\date{
Otoño 2025
} % Date, can be changed to a custom date

\newcommand{\pausa}{\pause} % Para usar una pausa en las presentaciones
%\newcommand{\pausa}{}      % Para que NO salgan las pausas

%------------------------------------------------------------
%The next block of commands puts the table of contents at the 
%beginning of each section and highlights the current section:
\AtBeginSection[]
{
  \begin{frame}
    \frametitle{Contenido}
    \tableofcontents[currentsection]
  \end{frame}
}
%------------------------------------------------------------

\begin{document}
{ % Portada de la universidad
  \setbeamertemplate{footline}{}
  \usebackgroundtemplate{\includegraphics[width=\paperwidth]{fondo-portada.png}}
  \begin{frame}
  \end{frame}
}

\begin{frame}
    \titlepage
\end{frame}

\usebackgroundtemplate{\includegraphics[width=\paperwidth]{fondo.png}}

% ex: ts=4 sw=4 sts=4 et filetype=tex
% SPDX-License-Identifier: CC-BY-SA-4.0
\begin{frame}
    \frametitle{Contenido}
    \tableofcontents
\end{frame}

\section{Algoritmo y sus características}

\begin{frame}[c]{Título}
    \begin{center}
        Texto
    \end{center}
\end{frame}

\section{Lenguajes de programación y su clasifiación}

\begin{frame}[c]{Título}
    \begin{center}
        Texto
    \end{center}
\end{frame}


{ % Portada de la universidad
  \setbeamertemplate{footline}{}
  \usebackgroundtemplate{\includegraphics[width=\paperwidth]{fondo-portada.png}}
  \begin{frame}
  \end{frame}
}
\end{document}
